% chapters/11/praktikum.tex
\begin{praktikumbox}
Praktikum ini adalah \textbf{Mini Proyek AI} pada RPS, yaitu pekerjaan
kelompok beranggotakan 2 sampai 3 mahasiswa. Keluarannya berupa kode program,
ringkasan desain, laporan \istilah{playtest} 2 sampai 3 halaman, dan presentasi
10 menit.

\begin{langkah}
  \item \textbf{Jalankan mini proyeknya.} Ketik ulang \code{ch11_mini_proyek.py}
        dari Daftar~\ref{lst:ch11-mini-proyek} ke dalam folder \code{code}.
        Berkas itu memakai \code{ch10_astar.py}, sehingga kedua
        berkas harus berada pada folder yang sama.

  \item \textbf{Jalankan evaluasi selang hitung ulang.} Jalankan program
        dengan pilihan \code{--evaluasi}, lalu catat tabel yang dihasilkan.
        Bandingkan dengan Tabel~\ref{tab:evaluasi-selang} pada buku ini, dan
        jelaskan perbedaan yang Anda temukan.
\begin{kodekeluaran}
python ch11_mini_proyek.py --evaluasi
\end{kodekeluaran}

  \item \textbf{Uji skenario pertama: ruang terbuka.} Jalankan permainan dan
        berlari lurus menjauh dari NPC. Catat jarak rata-rata yang ditampilkan
        dan kesan Anda mengenai perilaku NPC.

  \item \textbf{Uji skenario kedua: labirin.} Berlari mengelilingi lorong
        sehingga jalan lurus selalu terhalang. Catat banyaknya pencarian jalur
        dan kesan Anda mengenai jalur yang dipilih NPC.

  \item \textbf{Uji skenario ketiga: rintangan dinamis.} Tambahkan satu benda
        yang bergerak menutupi lorong, lalu amati perilaku NPC ketika jalurnya
        menjadi usang. Apabila perlu, perlambat laju bingkai atau perbesar
        selang hitung ulang agar kelemahannya lebih terlihat.

  \item \textbf{Uji satu perbaikan.} Pilih satu perbaikan dari
        Tabel~\ref{tab:perbaikan}---misalnya menyebar waktu perhitungan antar
        NPC---lalu terapkan dan ukur hasilnya dengan cara yang tercantum pada
        tabel itu.

  \item \textbf{Susun laporan \istilah{playtest}} yang memuat ketiga skenario
        beserta angka pengukuran, kesan mutu perilaku, perbaikan yang
        diterapkan beserta hasil pengukurannya, dan kendala yang dihadapi
        kelompok.

  \item \textbf{Siapkan presentasi 10 menit} yang memperagakan mini proyeknya
        secara langsung, lalu sampaikan pada pertemuan kelas. Setiap anggota
        kelompok harus memahami bagiannya masing-masing.

  \item \textbf{Simpan pekerjaan} ke repositori Git dengan pesan penyimpanan
        yang menerangkan isinya.
\end{langkah}
\end{praktikumbox}

\begin{table}[htbp]
\centering
\caption{Lembar kerja Mini Proyek AI}
\label{tab:lembar-mini-proyek}
{\setlength{\tabcolsep}{4pt}\footnotesize%
\begin{tabularx}{\textwidth}{@{}L{2.8cm} C{2.2cm} C{2.2cm} Y@{}}
\toprule
\headrow \textbf{Skenario} & \textbf{Jarak rata-rata} & \textbf{Banyaknya pencarian} & \textbf{Kesan mutu perilaku} \\
\midrule
Ruang terbuka & & & \\[1.6em]
\addlinespace
Labirin & & & \\[1.6em]
\addlinespace
Rintangan dinamis & & & \\[1.6em]
\addlinespace
\multicolumn{4}{@{}p{\textwidth}@{}}{\footnotesize\itshape
Perbaikan yang diterapkan: \dotfill} \\[1.4em]
\addlinespace
\multicolumn{4}{@{}p{\textwidth}@{}}{\footnotesize\itshape
Hasil pengukuran sesudah perbaikan: \dotfill} \\[1.4em]
\addlinespace
\multicolumn{4}{@{}p{\textwidth}@{}}{\footnotesize\itshape
Kendala yang dihadapi kelompok dan cara mengatasinya: \dotfill} \\[2.0em]
\bottomrule
\end{tabularx}}
\end{table}

\begin{catatabox}
Penilaian mini proyek ini memakai rubrik pada Lampiran~A: implementasi mesin
keadaan 25 persen, implementasi pencarian jalur 25 persen, penggabungan
keduanya 20 persen, pengujian beserta laporan 20 persen, dan presentasi
10 persen. Separuh nilainya berkaitan dengan pengujian dan
penggabungan, bukan dengan penulisan kode semata.
\end{catatabox}
